\section{CUDA GPU并行编程：矩阵乘法中的分块技术}

\subsection{引言}

在矩阵乘法的GPU优化中，\textbf{分块（Tiling）}是提升性能的核心技术。本章将深入探讨如何为矩阵乘法应用选
择最佳的分块策略，重点分析矩阵尺寸、SM 数量与分块尺寸三者之间的制约关系，特别是针对非方阵的处理方法。

\subsection{关键参数与约束}

在选择分块策略时，必须综合考虑以下三个关键参数：

\begin{enumerate}
    \item \textbf{矩阵尺寸（Matrix Dimensions）}：包括方阵与非方阵。对于非方阵，需在X和Y维度上采用不同的分块尺寸（如 $32 \times 16$），前提是矩阵维度必须是分块尺寸的整数倍。
    \item \textbf{SM数量（Number of SMs）}：直接决定GPU利用率。若分块总数少于SM数量，部分SM将处于空闲状态；若分块总数不是SM数量的整数倍，最后一轮计算将导致SM资源浪费。
    \item \textbf{分块尺寸（Tile Size）}：直接影响共享内存占用与执行时间。通常需要通过实验（Profiling）确定最优值，而非盲目选择。
\end{enumerate}

\noindent \textbf{核心原则}：为了最大化GPU占用率（Occupancy），\textbf{总分块数应尽可能为SM数量的整数倍}，且矩阵维度应能被分块尺寸整除。

\subsection{分块策略案例分析}

以下案例假设GPU拥有 \textbf{16个SM}，旨在说明不同参数组合对性能的影响。

\subsection{案例1：理想方阵分块}
\begin{itemize}
    \item \textbf{矩阵尺寸}：$40 \times 40$
    \item \textbf{分块尺寸}：$4 \times 4$
    \item \textbf{分析}：
    \[    \text{X方向分块数} = \frac{40}{4} = 10, \quad     \text{Y方向分块数} = \frac{40}{4} = 10    \]    \[    \text{总分块数} = 10 \times 10 = 100    \]    虽然100不是16的倍数，但 $100 \div 16 = 6 \dots4$，最后一轮仅有4个分块，利用率尚可。
    
    \item \textbf{更优解}：若选择分块尺寸使总分块数恰好为16或32等SM倍数，则利用率达到100\%。例如，若矩阵为 $64 \times 64$，分块为 $16 \times 16$，总分块数 $4 \times 4 = 16$，完美匹配16个SM。
\end{itemize}

\subsection{案例2：非方阵与边界填充}
\begin{itemize}
    \item \textbf{矩阵尺寸}：$31 \times N$
    \item \textbf{问题}：31无法被常见分块尺寸（如16、32）整除。
    \item \textbf{解决方案}：\textbf{零填充（Zero Padding）}。将矩阵扩展至 $32 \times N$，新增列/行补零。这虽引入少量冗余计算，但保证了内存访问的对齐与分块的规整性，整体收益远大于代价。
\end{itemize}

\subsection{案例3：分块数与SM不匹配}
\begin{itemize}
    \item \textbf{矩阵尺寸}：$40 \times 50$，\textbf{分块尺寸}：$10 \times 10$
    \item \textbf{分析}：
    \[    \text{总分块数} = \left(\frac{40}{10}\right) \times \left(\frac{50}{10}\right) = 4 \times 5 = 20    \]    20不是16的倍数。第一轮16个SM满载，第二轮仅剩4个分块，导致12个SM空闲，利用率骤降至25\%。
    \item \textbf{教训}：即使矩阵维度可被整除，若总分块数不是SM倍数，仍会造成严重的\textbf{尾部效应（Tail Effect）}。
\end{itemize}

\subsection{案例4：最差场景}
\begin{itemize}
    \item \textbf{矩阵尺寸}：$40 \times 41$，\textbf{分块尺寸}：$10 \times 10$
    \item \textbf{问题叠加}：
    \begin{enumerate}
        \item 41无法被10整除，产生不规则边界块（仅1列有效数据）。
        \item 总分块数 $4 \times 5 = 20$，非SM倍数，SM利用率低。
    \end{enumerate}
    此场景同时存在\textbf{内存访问不对齐}与\textbf{计算资源浪费}双重问题，应极力避免。
\end{itemize}

\subsection{非方形分块的灵活性}

分块不必是正方形。对于 $M \times K$ 与 $K \times N$ 的非方阵，可采用矩形分块以适配维度：
\begin{itemize}
    \item 示例：$64 \times 32$ 矩阵，可采用 $16 \times 8$ 分块。
    \item X方向：$64 / 16 = 4$ 块
    \item Y方向：$32 / 8 = 4$ 块
    \item 总分块数：$4 \times 4 = 16$，完美匹配16 SM。
\end{itemize}
\noindent \textbf{关键}：只要满足“维度可整除”且“总分块数为SM倍数”两个条件，任意矩形分块均为合法且潜在最优的选择。

\subsection{学习资源推荐}

建议关注 \textbf{NVIDIA GTC (GPU Technology Conference)}大会。该会议每年
3月举办，涵盖科学计算、AI、大模型、Omniverse及光线追踪等数百场专题讲座。可通过搜索引擎访问\textbf{GTC On-Demand}平台，免费注册后即可观看历年讲座文章与资料，是深入学习GPU并行编程的权威资源。

\subsection{总结与展望}

本章明确了分块策略选择的黄金法则：
\begin{center}
    \fbox{\parbox{0.9\textwidth}{\centering
        \textbf{最优分块} $\Longleftrightarrow$
        $\begin{cases}
            \text{矩阵维度} \mod \text{分块尺寸} = 0 \\
            \text{总分块数} \mod \text{SM数量} = 0
        \end{cases}$
    }}
\end{center}

下一章将转向 \textbf{Tensor Cores（张量核心）}。在Ampere及更新架构中，Tensor Core 执
行矩阵乘法的性能可达普通 CUDA Core 的\textbf{9倍以上}，已成为现代深度学习与高性能计算中GEMM操作的唯
一推荐路径。理解普通核心的分块原理，是掌握 Tensor Core 编程的重要基础。

